% foundations/ontology.tex

\chapter{Conceptual Structure}
\label{chap:conceptual-structure}
\index{Conceptual structure}

\section{The Ambiguity a Data Model Cannot Resolve Silently}

Suppose the edited CAD curve and a survey polyline differ by several
centimetres. A data model can store both coordinate sets, attach identifiers,
and calculate a residual. It cannot infer from those operations alone what the
difference means. The survey may observe the physical track; the CAD curve may
express an intended state; both may represent the same alignment; and the
residual may or may not be acceptable under the applicable engineering
knowledge.

Treating every record as another version of ``the alignment'' hides precisely
the distinctions needed to reason about the discrepancy. Treating every
different geometry as a different object loses continuity. Treating the measured
record as reality confuses evidence with its subject. Treating a computed pass
or fail as a decision transfers authority to an operation that does not own it.

The minimum remedy is to assign different formal roles to object, state,
realization, representation, observation, knowledge, evaluation, and decision.
Once separated, they can be related explicitly: an observation can concern a
physical state; a representation can expose an intended state; an evaluation
can compare them under applicable knowledge; a decision can use the resulting
evidence without becoming identical to it.

\section{The Engineering Object and Its States}

An engineering object supplies durable distinguishability across changes of
representation, realization, workflow, and project context. An alignment is
therefore not identical to a CAD curve, GIS feature, exchange entity, database
row, or rendered line. Multiple descriptions may concern the same alignment,
and a changed description does not by itself decide whether the engineering
object was modified or replaced.

A state describes that object under a specified engineering context. Intended,
metric, physical, operational, and evaluated states answer different questions
about one object. State change can therefore be discussed without making object
identity a side effect of coordinate equality.

For the alignment, intrinsic longitudinal order and curvature evolution provide
the constructive basis from which selected geometric states are derived. Later
chapters formalize those derivations; the conceptual requirement here is that
the derived state does not become the owner of identity.

\section{Realization and Representation}

Realization and representation answer different questions.

Metric realization asks how an intrinsic construction acquires measurable
spatial meaning under an explicit metric space, context, and applicable rules.
Physical realization asks how engineering intent becomes physical
infrastructure. Neither is merely a change of file format.

A representation selects information so that it can be communicated, stored,
exchanged, visualized, or processed. CAD geometry, GIS datasets, IFC entities,
sampled polylines, plots, and runtime structures are representations. A
representation may expose a realized state, but it cannot create missing
realization semantics or establish identity by itself.

This separation makes a useful question possible: does a difference arise from
the object's intended construction, from the chosen metric context, from
physical realization, from observation, or only from representation?

\section{Observation and Knowledge}

An observation is source-grounded information about an object, state, or
context, together with the provenance, method, time, conditions, and uncertainty
needed to interpret it as evidence. A survey polyline is therefore not simply
``the real alignment.'' It is an observation whose relation to a physical state
must be established.

Engineering knowledge is information admitted for engineering use within an
explicit scope, provenance, applicability, and authority. Observations may
support knowledge; storage or repetition does not admit them automatically.
Standards, engineering principles, derived results, and accountable decisions
may also support knowledge without becoming interchangeable sources.

The distinction allows conflicting observations to remain recorded while an
engineer determines what may be relied upon for the present task.

\section{Evaluation and Decision}

Evaluation applies knowledge, constraints, preferences, and comparison methods
to states or alternatives. It can derive a residual, identify a violation, rank
candidates, or show how the curvature edit propagates. These results are
evidence within an engineering process.

Decision is the accountable act that accepts, rejects, or selects within a
defined authority and context. A solver, viewer, or rule engine may support that
act but does not exercise it merely by producing output. This boundary is what
allows computation to become more capable without making responsibility less
clear.

\section{Relations Made Available}

With the roles separated, the thesis can state dependencies without category
errors:

\begin{itemize}
\item one engineering object may have several states and representations;
\item intrinsic construction may be metrically realized without transferring
  identity to coordinates;
\item a physical state may be observed without the observation becoming that
  state;
\item an observation may support knowledge through explicit interpretation and
  admission;
\item an evaluation may use applicable knowledge and produce evidence for a
  decision without becoming the decision.
\end{itemize}

These relations turn the earlier discrepancy into a tractable engineering
question. The intended and observed records can be related to one alignment,
their states and provenance can be identified, and a comparison can be made
under stated knowledge and context.

\section{Canonical Interpretation}

\begin{proposition}
Engineering objects, states, realizations, representations, observations,
engineering knowledge, evaluations, and decisions have distinct roles. Their
explicit relations preserve continuity while making selected consequences
computable.
\end{proposition}

This conceptual structure now licenses the formal development. The following
chapters state the axioms, notation, state models, and operators required to
express these relations precisely. They may formalize a dependency; they may not
erase the responsibility boundary that made the dependency meaningful.
